\documentclass{article}

\usepackage{amsmath,amssymb}
\usepackage{xurl}
\usepackage[hidelinks]{hyperref}
\setlength{\emergencystretch}{2em}

% Shared commands for notes in docs/paper-gaps/.

\DeclareUnicodeCharacter{2080}{\ifmmode{}_{0}\else\textsubscript{0}\fi}
\DeclareUnicodeCharacter{2081}{\ifmmode{}_{1}\else\textsubscript{1}\fi}
\DeclareUnicodeCharacter{2082}{\ifmmode{}_{2}\else\textsubscript{2}\fi}
\DeclareUnicodeCharacter{2099}{\ifmmode{}_{n}\else\textsubscript{n}\fi}

\newcommand{\Lean}{\textsc{Lean}}
\ExplSyntaxOn
\NewDocumentCommand{\leanid}{m}{
  \ifmmode
    \textsf{\detokenize{#1}}
  \else
    \group_begin:
    \urlstyle{sf}
    \tl_set:Nn \l_tmpa_tl {#1}
    \tl_replace_all:Nnn \l_tmpa_tl {\_} {_}
    \exp_args:NV \path \l_tmpa_tl
    \group_end:
  \fi
}
\ExplSyntaxOff
\providecommand{\tr}{\operatorname{tr}}
\newcommand{\tnleanIssueBase}{https://github.com/LionSR/TNLean/issues/}
\newcommand{\tnleanPrBase}{https://github.com/LionSR/TNLean/pull/}
\newcommand{\tnleanDocBase}{https://sirui-lu.com/TNLean/docs/}
\newcommand{\ghissue}[1]{\href{\tnleanIssueBase#1}{GitHub issue \##1}}
\newcommand{\ghpr}[1]{\href{\tnleanPrBase#1}{PR \##1}}
\newcommand{\doclink}[1]{\href{\tnleanDocBase#1.html}{\path{#1}}}

% Dirac notation and the identity operator, as in blueprint/src/macros/common.tex.
% \providecommand keeps note-local definitions authoritative.
\usepackage{dsfont}
\providecommand{\ket}[1]{|#1\rangle}
\providecommand{\bra}[1]{\langle #1|}
\providecommand{\braket}[2]{\langle #1 | #2 \rangle}
\providecommand{\ketbra}[2]{|#1\rangle\!\langle #2|}
\providecommand{\Id}{\mathds{1}}
\providecommand{\one}{\mathds{1}}
\providecommand{\C}{\mathbb{C}}
\providecommand{\MN}[1]{M_{#1}(\C)}
\providecommand{\MD}{M_D(\C)}

% External referencing for paper sources.  The build helper writes sanitized
% auxiliary files under build/paper-aux/ and excludes duplicated source labels.
\usepackage{xr}

% Import a generated paper auxiliary file with a caller-supplied label prefix.
% The candidate paths support builds from docs/paper-gaps/, the repository root,
% and build/paper-aux/ itself.
\newcommand{\importpaperaux}[2]{%
  \IfFileExists{../../build/paper-aux/#2.aux}{%
    \externaldocument[#1]{../../build/paper-aux/#2}%
  }{%
    \IfFileExists{build/paper-aux/#2.aux}{%
      \externaldocument[#1]{build/paper-aux/#2}%
    }{%
      \IfFileExists{#2.aux}{%
        \externaldocument[#1]{#2}%
      }{}%
    }%
  }%
}

% General external reference macros
\newcommand{\paperref}[2]{\ref{#1-#2}}
\newcommand{\papereqref}[2]{(\ref{#1-#2})}

% CPSV16 source (1606.00608 / MPDO-22-12-17-2.tex) external document and shortcuts
\importpaperaux{cpsv16-}{1606.00608/MPDO-22-12-17-2-tnlean}
\newcommand{\cpsvref}[1]{\paperref{cpsv16}{#1}}
\newcommand{\cpsveqref}[1]{\papereqref{cpsv16}{#1}}

% Verdict marker: \gapnote{<kind>}{<status>}. Kind is the policy
% classification (clarification, local-correction, scope-restriction,
% unfaithful, false-source, open-gap); status is open, wip (elimination
% underway), resolved, or historical. Machine-read by the site index;
% prints a verdict line.
\newcommand{\gapnote}[2]{\par\noindent\textbf{Verdict:} #1 (#2).\par}


\title{The Finite-Range Cyclic Knabe Coefficient}
\author{}
\date{\today}

\begin{document}
\maketitle

\section{Source statements}

For a nearest-neighbor chain, let $m$ be the number of local interaction
terms in each window and let $\gamma$ be the gap of the corresponding open
window Hamiltonian. Knabe's finite-size criterion gives the threshold
$\gamma>1/m$. This threshold originates in \cite{Knabe1988Energy}. It is
quoted for MPS parent Hamiltonians in lines~1483--1489 of the local source for
\cite{PerezGarcia2007Matrix} and in lines~2194--2197 of the local source for
\cite{Cirac2021Matrix}. The quoted passages state only the threshold. The
closed-form periodic coefficient
\begin{align}
    \delta&=\frac{m\gamma-1}{m-1}
\end{align}
is obtained by the present derivation, as the range-two case of the coefficient
below.

\section{Finite-range derivation}

Let $h_i$ be cyclically indexed orthogonal projections, let
$H=\sum_i h_i$, and define
\begin{align}
    A_s&=\sum_{a=0}^{m-1}h_{s+a},
    &Q_q&=\sum_i\bigl(h_i h_{i+q}+h_i h_{i-q}\bigr).
\end{align}
Assume $1\leq R\leq m$, $N\geq2m$, and that terms at cyclic separation at
least $R$ commute. Cyclic double counting gives
\begin{align}
    \sum_s A_s^2&=mH+\sum_{q=1}^{m-1}(m-q)Q_q.
\end{align}
The projection inequality $h_i h_j+h_j h_i\leq h_i+h_j$ gives $Q_q\leq2H$.
The separated terms satisfy $Q_q\geq0$ for $q\geq R$, and expansion of $H^2$
gives
\begin{align}
    H+\sum_{q=1}^{m-1}Q_q&\leq H^2.
\end{align}
Set $c=m-R+1$, so that $m-q=c+(R-1-q)$. The correction term splits as
\begin{align}
    \sum_{q=1}^{m-1}(R-1-q)Q_q
    &=\sum_{q=1}^{R-2}(R-1-q)Q_q
      +\sum_{q=R}^{m-1}(R-1-q)Q_q,\\
    \sum_{q=1}^{R-2}(R-1-q)Q_q
    &\leq 2H\sum_{q=1}^{R-2}(R-1-q),\\
    \sum_{q=R}^{m-1}(R-1-q)Q_q
    &\leq0.
\end{align}
Indeed, the first sum has positive weights and $Q_q\leq2H$ termwise. In the
second sum, $R-1-q\leq0$ and $Q_q\geq0$: the latter estimate applies because
$R\leq q\leq m-1\leq N-R$. The triangular sum is
\begin{align}
    \sum_{q=1}^{R-2}(R-1-q)
    &=\frac{(R-1)(R-2)}{2}.
\end{align}
Therefore
\begin{align}
    \sum_s A_s^2
    &\leq cH^2+\bigl(m-c+(R-1)(R-2)\bigr)H,\\
    m-c+(R-1)(R-2)
    &=(R-1)^2,\\
    \sum_s A_s^2
    &\leq cH^2+(R-1)^2H.
\end{align}
Consequently, the local inequalities $A_s^2\geq\gamma A_s$ imply
\begin{align}
    H^2&\geq\frac{m\gamma-(R-1)^2}{m-R+1}H.
\end{align}
The condition $(R-1)^2<m\gamma$ makes this coefficient positive. At $R=2$
it reduces to Knabe's nearest-neighbor coefficient.

\section{Verdict}

The nearest-neighbor threshold is source-attributed as above. The coefficient
for arbitrary finite range is a TNLean derivation, formalized by the cyclic
projection theorem.\footnote{Formal declaration:
\leanid{ProjectionGeometry.quadraticForm_sum_projections_of_cyclic_knabe}.}
It is not a statement attributed
to Knabe, to the gap-from-hypotheses criterion
\cite[Theorem~2.1(i)]{Nachtergaele1996Spectral} of Nachtergaele's martingale
argument, or to the two MPS reviews. The abstract result concerns a cyclic family of
projections. Applying it to every sufficiently large MPS parent Hamiltonian
still requires a uniform gap for the corresponding nonwrapping open
interaction.

\bibliographystyle{alpha}
\bibliography{references}

\end{document}
